\documentclass[11pt,a4paper]{amsart}

\usepackage[utf8]{inputenc}
\usepackage[T1]{fontenc}
\usepackage[ngerman]{babel}
\usepackage{amsmath,amssymb,amsthm}
\usepackage{array,booktabs,longtable}
\usepackage{hyperref}

\newtheorem{theorem}{Satz}[section]
\newtheorem{proposition}[theorem]{Proposition}
\newtheorem{remark}[theorem]{Bemerkung}
\newtheorem{openproblem}[theorem]{Offenes Problem}

\newcommand{\LFF}{\mathrm{LFF}}
\newcommand{\OPEN}{\mathrm{OPEN}}

\title[P10 --- No-Go Theorems]{P10 --- No-Go Theorems for Canonical Global Coupling}
\author{Objekt-X-Programm}
\date{2026-08-09}

\begin{document}

\begin{abstract}
P10 sammelt die bis zum 9. August 2026 eindeutig auditierten Negativbefunde aus den
SYN-Blöcken P05--P09 und ordnet sie nach mathematischem Typ. Ziel ist nicht,
den Suchraum künstlich zu verkleinern, sondern falsche Abkürzungen dauerhaft
zu sperren. Lokale Primfaser- und Rang-eins-Modelle liefern nicht automatisch
die globale positive Geometrie von Objekt~X; der auditierte
Primfaser-Transportgenerator ist kein diskreter Hilbert--Pólya-Endoperator;
eine endliche Feshbachidentität beweist keine globale Schatten-/Fredholm-
Grenzstruktur; und der konkrete NEU-088--90-Determinantenpfad mit
$h_r=r$ und $M_N=N/\log N$ kollabiert auf $D_N(z)\to1$ statt auf einen
nichttrivialen $C\xi(z)$-Grenzwert. Zugleich bleibt
$\LFF\Rightarrow\mathrm{Rampe}$ nur einseitig bewiesen und
$\mathrm{Rampe}\Rightarrow\LFF$ offen. Mehrere historische Jacobi-,
Primeclock-, Mellin-, Finite-Part-, KMS-, zyklische und Hopf-Reparaturen
scheitern in ihrem konkreten Scope. Keiner dieser Befunde schließt eine
globale nichtorthogonale positive Gram-/Hilbertraumstruktur mit
Primzahlpotenzkanälen, archimedischem Anteil und echter globaler Kopplung aus.
P10 enthält keinen RH-Beweis und konstruiert Objekt~X nicht.
\end{abstract}

\maketitle

% ----------------------------------------------------------------
\section{Statussemantik und Anti-Overreach-Firewall}
% ----------------------------------------------------------------

\begin{remark}[Vier logisch verschiedene Klassen]\label{rem:status}
P10 trennt strikt:
\begin{enumerate}
  \item \emph{NO-GO}: eine exakt formulierte Behauptung oder Konstruktion ist
        im angegebenen Scope widerlegt;
  \item \emph{Kandidaten-/Scope-No-Go}: ein konkreter Modellpfad oder eine
        Schlussweise scheitert, ohne Alternativen auszuschließen;
  \item \emph{SUPERSEDED}: eine historische Formel oder Typisierung ist durch
        einen späteren korrigierten Stand ersetzt;
  \item \emph{OPEN / CONDITIONAL}: weder bewiesen noch widerlegt und damit
        ausdrücklich kein No-Go.
\end{enumerate}
\end{remark}

\begin{remark}[Präzedenzregel]\label{rem:precedence}
Verbindlich ist die Reihenfolge
\[
\text{P10 FINAL SEAL / Matrix}
>
\text{P10 Gegencheck / Targeted-Reaudit}
>
\text{lokal synchronisierter SYN-Stand}
>
\text{historischer NEU-Stand}.
\]
\end{remark}

\begin{remark}[Retired Slot P10-N15]\label{rem:n15}
Der historische Slot P10-N15 ist kein aktiver No-Go. Der auditierte P07-Block
beweist
\[
  \boxed{\LFF\Longrightarrow\mathrm{Rampe}},
\]
aber für
\[
  \boxed{\mathrm{Rampe}\Longrightarrow\LFF}
\]
liegt weder Beweis noch Gegenbeispiel vor. Der Slot ist daher
\emph{RETIRED / MOVED TO P10-O29}.
\end{remark}

% ----------------------------------------------------------------
\section{Primkanal-, Lift-, Projektor- und Spektralfirewalls}
% ----------------------------------------------------------------

\begin{theorem}[Gewichteter Rang-eins-Kanal]\label{thm:n01}
Für das auditierte induzierte Rang-eins-Modell
\[
  P_p=|c_p|^2\Pi_p^{(1)}
\]
gilt
\[
  P_p^2=|c_p|^2P_p.
\]
Ohne zusätzliche Normierung $|c_p|^2=1$ ist $P_p$ daher nicht automatisch ein
orthogonaler Projektor. Andere normierte Projektoren und andere
Primkanalrealisierungen bleiben möglich. \emph{[P10-N01]}
\end{theorem}

\begin{theorem}[Endlich viele Primlabels erzwingen keinen endlichen Rang]
\label{thm:n02}
Aus endlich vielen Primlabels bis $N$ folgt nicht
\[
  \operatorname{rank}K_N\le\pi(N),
\]
wenn die Primfasern weiterhin unendlich viele interne Indizes besitzen.
Eine zusätzliche echte Fasertrunkierung kann den Rang dagegen endlich machen.
\emph{[P10-N02]}
\end{theorem}

\begin{theorem}[Source-Cone-Firewall]\label{thm:n03}
Im explizit auditierten Source-Cone liefern zusätzliche nichttriviale lineare
Abbildungen $L_{p,a}$ keinen neuen homogenen Kern der dort verlangten Art.
Dieser Befund betrifft ausschließlich die geprüfte Quellarchitektur;
Operatoren außerhalb dieses Source-Cones und andere intrinsische
Quellkonstruktionen bleiben offen. \emph{[P10-N03]}
\end{theorem}

\begin{theorem}[Primfaser-Transport ist kein diskreter HP-Endoperator]
\label{thm:n04}
Der auditierte Transportgenerator vom Typ
\[
  2i\kappa\frac{d}{dt}
\]
hat auf der geprüften Primfaser rein absolut kontinuierliches Spektrum und
keinen kompakten Resolventen. Er ist in dieser Realisierung kein diskreter
Hilbert--Pólya-Endoperator. Ein späterer global gekoppelter Endoperator, eine
andere Hilbertisierung oder zusätzliche zusammengesetzte Sektoren werden nicht
ausgeschlossen. \emph{[P10-N04; N05 SUPERSEDED]}
\end{theorem}

\begin{theorem}[Matrixkoeffizient ist kein automatisches Normquadrat]
\label{thm:n06}
Aus der Unitarität von $U_t$ folgt für
\[
  g_a(\log p)=\operatorname{Re}\langle a,U_{\log p}a\rangle
\]
nicht automatisch $g_a(\log p)\ge0$. \emph{[P10-N06]}
\end{theorem}

\begin{theorem}[Naive Mangoldt-Energielesart]\label{thm:n07}
Eine historische gradnormierte Prime-edge-Energie reproduziert auf gemischten
Zahlen nicht automatisch die Mangoldt-Funktion; insbesondere ist
\[
  \Lambda(6)=0.
\]
Ein echter Prime-Power-/Mangoldt-Mediator bleibt möglich. \emph{[P10-N07]}
\end{theorem}

\begin{proposition}[Historische Typ- und Symmetrieformeln]\label{prop:n08n10}
Im auditierten Jacobi-/Divisorgraph-Scope gelten:
\begin{enumerate}
  \item $J_N^-:=\frac12(\Theta_N-\Theta_N^\dagger)$ ist schiefadjungiert;
        $S_N=-iJ_N^-$ ist der selbstadjungierte Kandidat. \emph{[P10-N08]}
  \item Für $k>1$ gilt nicht $\log(p^k)=\Lambda(p^k)$. \emph{[P10-N09]}
  \item Eine reine $r$-Gradierung erzwingt ohne zusätzliche echte
        Bipartitheit nicht das Verschwinden aller ungeraden Spuren.
        \emph{[P10-N10]}
\end{enumerate}
\end{proposition}

% ----------------------------------------------------------------
\section{Feshbach-, Fredholm- und Determinantenfirewalls}
% ----------------------------------------------------------------

\begin{theorem}[Finite Feshbachidentität ist keine globale Konvergenztheorie]
\label{thm:n11}
Eine endliche Feshbach-/Schur-Komplement-Identität beweist für sich allein
weder globale Operatorgrenzen noch Hilbert--Schmidt-, Spurklassen- oder
Fredholm-Konvergenz. Dafür sind separat bewiesene uniforme Schatten- und
Grenzabschätzungen erforderlich. \emph{[P10-N11]}
\end{theorem}

\begin{theorem}[Reconciliertes Determinanten-No-Go im NEU-088--90-Scope]
\label{thm:n12}
Im konkreten Modell
\[
  h_r=r,\qquad M_N=\frac{N}{\log N},\qquad z\text{ fest und zulässig},
\]
gilt nach dem späteren P06-Reaudit
\[
  T_N(z)=O_z\!\left(\frac{\log\log N}{\log N}\right)\to0,
\]
ferner
\[
  \|C_N(z)\|_{HS}\to0,
  \qquad
  \operatorname{Tr}(C_N(z)^k)\to0\quad(k\ge2),
\]
und damit
\[
  \boxed{D_N(z)\to1.}
\]
Diese konkrete Skalierung liefert keinen nichttrivialen
$C\xi(z)$-Determinantengrenzwert. \emph{[P10-N12]}
\end{theorem}

\begin{remark}[SUPERSEDED-Historie]\label{rem:n13n14}
Die historische Aussage
\[
  D_N(z)\to e^{-\gamma^2/4}
\]
ist für denselben Modellscope \emph{SUPERSEDED}; sie ist kein zweiter No-Go.
\emph{[P10-N14]} Für komplexes $z$ gilt allgemein
\[
  \|C\|_{HS}^2=\operatorname{Tr}(C^*C),
\]
nicht $\operatorname{Tr}(C^2)$. Die historische Selbstadjungiertheits-/
Hilbert--Schmidt-Lesart ist entsprechend ersetzt. \emph{[P10-N13]}
\end{remark}

\begin{remark}[Was der Determinantenbefund nicht ausschließt]
\label{rem:det-open}
Offen bleiben andere Skalierungen, zusätzliche Renormierungen, globale
Feshbach-Transfers, Fredholm- oder $\det_2$-Konstruktionen und eine
Determinantenstruktur nach vorheriger Weil-Hilbertisierung.
\end{remark}

% ----------------------------------------------------------------
\section{Weil-, Formfaktor- und Herglotzfirewalls}
% ----------------------------------------------------------------

\begin{theorem}[LFF ist keine vollständige Weilform-Konstruktion]
\label{thm:n16}
Der lokale Formfaktor allein konstruiert oder identifiziert $Q_{\rm Weil}$
nicht. Eine zusätzliche Typisierungs-, Autokorrelations- oder
Positivierungsbrücke bleibt erforderlich. \emph{[P10-N16]}
\end{theorem}

\begin{theorem}[Massendiskrepanz im unrenormierten Spektralmaß]
\label{thm:n17}
Für einen normierten zyklischen Vektor gilt
\[
  \mu_{\Omega,N}(\mathbb R)=1,
\]
während das arithmetische Zielmaß unendliche Gesamtmasse besitzt. Das
unskalierte Wahrscheinlichkeitsmaß genügt daher nicht als direkte
Approximation. Eine Renormierung $c_N>0$ und Nevanlinna-/Tail-Kontrolle bleiben
erforderlich. \emph{[P10-N17]}
\end{theorem}

\begin{theorem}[Vage Maßkonvergenz genügt nicht automatisch]
\label{thm:n18}
Aus $\widetilde\mu_N\to\mu_{\rm arith}$ im vagen Sinn folgt bei unendlicher
Zielmasse nicht automatisch lokal gleichmäßige Konvergenz der zugehörigen
Nevanlinna-/Herglotzfunktionen. Geeignete Tail- und Gewichtskontrolle muss
separat bewiesen werden. \emph{[P10-N18]}
\end{theorem}

\begin{remark}[SUPERSEDED Pol-Lesart]\label{rem:n19}
Pole $\pm i/2$ gehören nicht als Pole in das kanonisch zentrierte
Nullstellen-Herglotzobjekt
\[
  m_{\rm arith}=-\Xi'/\Xi.
\]
Gamma-/Polbeiträge gehören auf die getrennte explizite-Formel-Seite.
\emph{[P10-N19]}
\end{remark}

% ----------------------------------------------------------------
\section{Jacobi-, Prä-Lanczos- und Renormierungsfirewalls}
% ----------------------------------------------------------------

\begin{theorem}[Unrenormierter erster Lanczos-Kanal kollabiert]
\label{thm:n20}
Für den auditierten Startvektor $q_0=e_1$ gilt
\[
  b_{1,N}\sim\gamma\sqrt{\frac{\log N}{N}}\to0.
\]
Diese unrenormierte direkt symmetrisierte Jacobi-Folge besitzt daher keinen
nichtdegenerierten ersten Grenzparameter $b_1>0$. \emph{[P10-N20]}
\end{theorem}

\begin{theorem}[Startvektor-Resolvente liefert nicht $m_{\rm arith}$]
\label{thm:n21}
Im selben unrenormierten Kanal konvergiert die Startvektor-Resolvente gegen
\[
  -\frac1z,
\]
nicht gegen $m_{\rm arith}$. \emph{[P10-N21]}
\end{theorem}

\begin{theorem}[Erste Lanczos-Kante entscheidet nicht den Gesamtoperator]
\label{thm:n22}
Aus $b_{1,N}\to0$ folgt nicht, dass alle höheren Offdiagonalen verschwinden
oder ein möglicher Grenzoperator global diagonal ist. \emph{[P10-N22]}
\end{theorem}

\begin{proposition}[Self-Energy-, Skalierungs- und Schur-Firewalls]
\label{prop:n23n25}
Folgende historische Schritte sind gesperrt:
\begin{enumerate}
  \item Vektorwertige Self-Energy und skalare Quadratformsumme dürfen nicht
        identifiziert werden. \emph{[P10-N23, SUPERSEDED]}
  \item Aus $P_{kl}\le Cc^{1/2}$ und $A_{kl}=P_{kl}c^{1/2}$ folgt
        $A_{kl}=O(c)$, nicht $O(1)$. \emph{[P10-N24]}
  \item Phasige Cancellation kontrolliert nicht automatisch
        $\sup_i\sum_{j\ne i}|T_{ij}|$. \emph{[P10-N25]}
\end{enumerate}
\end{proposition}

\begin{theorem}[Ungewichtete Primeclock-H1-Schranke]
\label{thm:n26n27}
Die historische Behauptung
\[
  \left|\sum_{p\in[P,2P]}p^{-iu}\right|\le\frac{C}{|u|}
\]
mit $P$-unabhängigem $C$ ist für festes $u$ falsch. Damit ist auch der konkret
hiervon abhängige NEU-133-Abel/H1-Kern in dieser ungewichteten Form nicht
bewiesen. Ein korrekt gewichtetes Primeclock-/Abel-Ersatzlemma bleibt offen.
\emph{[P10-N26, P10-N27]}
\end{theorem}

\begin{proposition}[Spur- und Normierungsfirewalls]\label{prop:n28n32}
Es gelten die korrigierten Firewalls:
\begin{enumerate}
  \item
  \[
    \sum_{p\le N}\frac{(\log p)^2}{p}
    \sim\frac12(\log N)^2,
  \]
  nicht $\frac13(\log N)^3$. \emph{[P10-N28, SUPERSEDED]}
  \item Divergenz eines Rohanteils folgt nicht allein aus einer oberen
        Schranke für $|c_p|^2$. \emph{[P10-N29]}
  \item Allgemein gilt nicht $\|CC^*\|_{S_1}=\|C\|_{op}^2$; für
        Hilbert--Schmidt-$C$ gilt
        \[
          \|CC^*\|_{S_1}=\|C\|_{S_2}^2.
        \]
        \emph{[P10-N30]}
  \item Spurklasse allein liefert keine primeweise Eigenwert-/Eulerprodukt-
        Lesart ohne zusätzliche Primdiagonalität/T2. \emph{[P10-N31]}
  \item Die historische zweite-Spur-Formel mit zusätzlichem Gesamtfaktor ist
        ersetzt durch
        \[
          \operatorname{Tr}(\Sigma^2)
          =\sum_{p,q}w_pw_q\operatorname{Tr}(P_pP_q).
        \]
        \emph{[P10-N32, SUPERSEDED]}
\end{enumerate}
\end{proposition}

\begin{theorem}[$\vartheta$ ist nicht $\psi$: Prime-only-Mellin-Firewall]
\label{thm:n33n36}
Ein Prime-only-Cutoff darf die Nullstellenterme der vollen
$\psi$-Explizitformel nicht ohne Prime-Power-Korrektur übernehmen. Der direkte
Import von
\[
  -\frac{\zeta'}{\zeta}(\beta+s)
\]
für die historische Prime-only-Summe $S_{\varphi,X}$ ist nicht korrekt; der
passende volle Mangoldt-Kanal ist $\Psi_{\varphi,X}$.
\emph{[P10-N33, P10-N34]}

Für die dort verwendete Mellintransformierte gilt außerdem nicht allgemein
``ganz und $\widehat\varphi(0)=1$''; im auditierten Setup liegt ein einfacher
Pol mit Residuum $1$ vor. \emph{[P10-N35, SUPERSEDED]}

Die historische $\Psi-S$-Differenzformel ist ersetzt; die korrigierte
Prime-Power-Differenz lautet
\[
  \sum_{k\ge2,p}\log p\,
  \bigl[\varphi(p^k/X)-\varphi(p/X)\bigr]p^{-k\beta}.
\]
\emph{[P10-N36]}
\end{theorem}

\begin{theorem}[Analytische Fortsetzung ist noch keine operatorielle
Finite-Part-Realisierung]\label{thm:n37}
Eine Definition
\[
  \operatorname{Tr}_{\rm reg}
  :=\operatorname{AC}\!\left[-\frac{\zeta'}{\zeta}\right]
\]
beweist für sich allein keinen operatoriellen Cutoff-/Finite-Part-Grenzwert.
Eine echte Operatorrealisierung bleibt offen. \emph{[P10-N37]}
\end{theorem}

% ----------------------------------------------------------------
\section{Hochschild-, KMS-, zyklische und Hopf-Firewalls}
% ----------------------------------------------------------------

\begin{theorem}[Frühe geladene HH4-Schablone]\label{thm:n38}
Der frühe symmetrische geladene Grad-4-Kandidat besitzt im auditierten Scope
verschwindende volle Alternation. Das schließt diesen Kandidaten aus, nicht
beliebige geladene $HH^4$-Klassen. \emph{[P10-N38]}
\end{theorem}

\begin{theorem}[NEU-212-Schwartz-Zieltyp]\label{thm:n39n40}
Für die dort definierte absolute Schnellabfallbedingung gilt nicht
\[
  A_{\rm alg}\subset A^\infty.
\]
Schon $1$ und die algebraischen Erzeuger liegen nicht im behaupteten Zieltyp.
Auch die konkrete Regularisierung $G/\log(\nu+2)$ ist nicht Schwartz. Diese
historische Route ist durch logarithmische Zieltypen ersetzt.
\emph{[P10-N39, P10-N40]}
\end{theorem}

\begin{theorem}[Schwartz-Inkremente und divergierende Shellgewichte]
\label{thm:n41}
Die für die singuläre Potentialroute benötigten divergierenden Gewichte sind
mit der historischen Schwartz-Inkrementforderung im auditierten Aufbau nicht
vereinbar. Andere Regularitätsarchitekturen werden nicht ausgeschlossen.
\emph{[P10-N41]}
\end{theorem}

\begin{theorem}[Kein nichttrivialer universeller globaler Bimodul-Glätter]
\label{thm:n42}
Ein globaler normstetiger $A_{\rm alg}$-Bimoduloperator
\[
  R:A_{C^*}\to A^\infty\subsetneq A_{C^*}
\]
muss im auditierten unitalen Zentralisator-Scope trivial sein. Ein
nachträglicher universeller Glätter repariert daher den historischen Zieltyp
nicht. Direkt definierte logarithmische Zielräume bleiben möglich.
\emph{[P10-N42; P09-CORE-NOGO]}
\end{theorem}

\begin{theorem}[Reguläre Implementierer und konkrete L/R/S-Platzierungen]
\label{thm:n43n44}
Normkonvergente Potentialimplementierer liefern im relevanten Quotienten nur
innere/invisible Derivationen. Außerdem scheitern die drei konkret auditierten
dyadischen L/R/S-Platzierungen. Die relation-adaptierte Architecture~III bleibt
offen. \emph{[P10-N43, P10-N44]}
\end{theorem}

\begin{remark}[SUPERSEDED Quotientenroute]\label{rem:n45}
Der frühe Baker-/Log-Gewichts-Separationsansatz beweist den vollen Quotienten
$M/[A,M]$ nicht. Für den positiven Cup-Beweis genügt ein partieller Quotient;
der volle Quotient bleibt offen. \emph{[P10-N45, SUPERSEDED]}
\end{remark}

\begin{theorem}[Direkter geladener KMS-Detektor]\label{thm:n46}
Für ein homogenes nichtneutrales Zielelement und $\beta>0$ verschwindet die
direkte KMS-Auswertung. Nichtverschwindende Detektoren benötigen eine explizite
Gradneutralisierung. \emph{[P10-N46]}
\end{theorem}

\begin{theorem}[Konkreter I4-KMS-Kozykel]\label{thm:n47}
Im bewiesenen Bereich $\beta>1$ gilt
\[
  T_{\sigma_\beta}\Phi_{\beta,\chi}
  =g^{-\beta}\Phi_{\beta,\chi},
  \qquad g\ne1.
\]
Der konkrete I4-Repräsentant ist dort nicht bereits ein invariant
getwistet-zyklischer Repräsentant. Andere Repräsentanten oder
Koeffiziententheorien bleiben offen. \emph{[P10-N47]}
\end{theorem}

\begin{theorem}[Standard-Zyklisierung annihiliert nichttriviale
Gewichtssektoren]\label{thm:n48}
In einem parazyklischen Eigenraum mit Gewicht $w\ne1$ ist $1-T$ invertierbar;
der Sektor verschwindet daher unter gewöhnlicher Invarianten-/Koinvarianten-
Zyklisierung. \emph{[P10-N48]}
\end{theorem}

\begin{theorem}[Externe Eigenlinie ist noch keine zyklische
Koeffiziententheorie]\label{thm:n49}
Eine formale eindimensionale Eigenlinie kompensiert zwar einen Eigenwert,
liefert aber ohne passende Koflächen-, Kodegenerations- und Rotationsstruktur
keine zyklische Koeffiziententheorie. \emph{[P10-N49]}
\end{theorem}

\begin{theorem}[Zwei minimale Koeffizientenreparaturen scheitern]
\label{thm:n50n51}
Im auditierten Scope scheitern sowohl ein eindimensionales unital-
nichtdegeneriertes $\sigma_\beta$-äquivariantes $A_{\rm alg}$-Bimodul als auch
der minimale Standard-SAYD-Pfad über $H_\beta=\mathbb C[\mathbb Z]$, wenn
exakter KMS-Twist, Ladungskompensation und Stabilität gleichzeitig verlangt
werden. Ein nichtstandardmäßiger $A$-relativer Hopf-Koeffizient bleibt offen.
\emph{[P10-N50, P10-N51]}
\end{theorem}

\begin{theorem}[Unmarkierte Orbitmodule verlieren den Orbitindex]
\label{thm:n52}
Für die gesättigten unmarkierten Orbitmodule gilt
\[
  N_k=N_0.
\]
Sie können verschiedene Orbitgrade daher nicht injektiv kodieren. Eine extern
markierte Orbitsumme bleibt möglich. \emph{[P10-N52]}
\end{theorem}

\begin{theorem}[Kanonischer skalarer Basislift besitzt keine konstante globale
Rotationseigenrelation]\label{thm:n53}
Für den kanonischen Basislift $\widetilde L_0$ und den daraus gebildeten
Kozykel $\Phi_0$ gilt im bewiesenen KMS-Bereich $\beta>1$
\[
  \boxed{t\Phi_0\ne C\Phi_0\qquad\forall C\in\mathbb C.}
\]
Andere zyklische/getwistet-zyklische Repräsentanten, orbitverschiebende
nichtkanonische Lifts, andere Koeffizienten und Weil-/Gamma-Korrekturen bleiben
offen. \emph{[P10-N53; P09-CORE-NOGO]}
\end{theorem}

\begin{remark}[Historische Rotationsformeln]\label{rem:n54}
Die Formeln $t\Phi_0=g^{-\beta}\Phi_0$ und $s=-1$ für den kanonischen Basislift
sind \emph{SUPERSEDED}. \emph{[P10-N54]}
\end{remark}

% ----------------------------------------------------------------
\section{Positiver Restbestand}
% ----------------------------------------------------------------

\begin{remark}[P10 ist eine Firewall-Sammlung]
P05--P09 liefern zugleich positive Bausteine: relative Primkanäle und ihre
Quell-/Liftgeometrie, finite Feshbach-/Schur-Komplement-Identitäten im korrekten
endlichen Scope, den Weilform-/Herglotz-Rahmen, korrigierte Prime-Power-/
Mangoldt- und Spurformeln sowie im P09-Strang eine nichttriviale geladene
äußere Derivation in größerem Koeffizientenraum und den nichttrivialen geladenen
Hochschild-Cup in $\mathfrak M_{\rm glob}^{\log}$. P10 ist deshalb keine
Negativbilanz des Gesamtprogramms.
\end{remark}

% ----------------------------------------------------------------
\section{OPEN / CONDITIONAL: ausdrücklich außerhalb des No-Go-Scope}
% ----------------------------------------------------------------

\small
\begin{longtable}{>{\raggedright\arraybackslash}p{0.12\textwidth}
                  >{\raggedright\arraybackslash}p{0.62\textwidth}
                  >{\raggedright\arraybackslash}p{0.18\textwidth}}
\toprule
ID & Offener Punkt & Status\\
\midrule
\endfirsthead
\toprule
ID & Offener Punkt & Status\\
\midrule
\endhead
P10-O01 & $c_p\ne0$ für alle Primkanäle & OPEN\\
P10-O02 & Liftunabhängigkeit und universelle Asymptotik von $|c_p|^2$ & OPEN\\
P10-O03 & neuer intrinsischer Ursprung von $L_3$ / vollständiges Zieltuple & OPEN\\
P10-O04 & voller balancierter Prime-Power-Lift $h_n^{\rm bal}=n^{-1/2}I$ & OPEN\\
P10-O05 & globale Primorthogonalität bzw. globale Kreuzblöcke & OPEN\\
P10-O06 & intrinsische $\gamma_N=1$-Rigidität & OPEN\\
P10-O07 & $S_4\setminus S_2$-Grenzstruktur / globale Schatten-Fredholm-Brücke & OPEN\\
P10-O08 & Selbstadjungiertheit des historischen $A_N^{\rm Jac,-}$ & OPEN\\
P10-O09 & kanonische Nevanlinna-Renormierung $(c_N,a_N)$ und Tail-Kontrolle & OPEN\\
P10-O10 & $m_{\rm arith}=\Pi_\gamma(X)$ & OPEN\\
P10-O11 & $b_{2,N}/b_{1,N}\to\infty$ & OPEN\\
P10-O12 & allgemeines skalares Renormierungs-No-Go & CONDITIONAL auf O11\\
P10-O13 & intrinsische positive nichtskalare Prä-Lanczos-Metrik $W_N$ & OPEN\\
P10-O14 & intrinsische Schranke $|c_p|^2=O((\log p)^2/p)$ & OPEN / CONDITIONAL\\
P10-O15 & intrinsisches T2 und $c_p\ne0$ & OPEN\\
P10-O16 & gewichtetes Primeclock-/Abel-Ersatzlemma & OPEN\\
P10-O17 & quantitativer/uniformer $\Psi/S$-Transfer & OPEN\\
P10-O18 & uniforme nullstellenvermeidende Kontur + volle Residuenzählung & OPEN\\
P10-O19 & operatorielle $\operatorname{Tr}_{\rm reg}$-/Finite-Part-Realisierung & OPEN\\
P10-O20 & lokaler $M_{g,p}^{\log}$ als voller $A_{(p),\rm alg}$-Bimodul & OPEN\\
P10-O21 & voller Quotient $M/[A,M]$ & OPEN\\
P10-O22 & $HH^1(A_{\rm alg},A_{\rm alg})_g$ und $HH^4(A_{\rm alg},A_{\rm alg})_g$ & OPEN\\
P10-O23 & $\beta=1$ in der I4-Gibbs-Auswertung & OPEN\\
P10-O24 & anderer zyklischer/getwistet-zyklischer Repräsentant & OPEN\\
P10-O25 & genuin orbitverschiebender nichtkanonischer Lift & OPEN\\
P10-O26 & nichtstandardmäßiger $A$-relativer Hopf-Koeffizient & OPEN\\
P10-O27 & NEU-205 Architecture III & OPEN\\
P10-O28 & Weil-/Gamma-Korrektur des zyklischen/kohomologischen Pfads & OPEN\\
P10-O29 & Rampenform $\Rightarrow$ LFF & OPEN\\
\bottomrule
\end{longtable}
\normalsize

% ----------------------------------------------------------------
\section{Objekt-X-Firewall}
% ----------------------------------------------------------------

\begin{theorem}[P10 ist kein Existenz-No-Go für Objekt X]\label{thm:objectx-firewall}
Aus P10 folgt nicht
\[
  \text{``Objekt X existiert nicht''.}
\]
Außerhalb der ausgeschlossenen konkreten Pfade bleiben insbesondere
\[
  \boxed{
  \text{globale nichtorthogonale Gramkopplung}
  +\text{ Primzahlpotenzkanäle}
  +\text{ archimedischer Kanal}
  +\text{ positive globale Weil-Geometrie}.}
\]
Die offene Hauptfrage bleibt, ob eine kanonische globale Quelle und Kopplung
existiert, deren positive Vervollständigung die vollständige Weilform
realisiert.
\end{theorem}

% ----------------------------------------------------------------
\section{Kompakte Statusmatrix}
% ----------------------------------------------------------------

\small
\begin{longtable}{>{\raggedright\arraybackslash}p{0.29\textwidth}
                  >{\raggedright\arraybackslash}p{0.62\textwidth}}
\toprule
Familie & Bindender P10-Endstand\\
\midrule
\endfirsthead
\toprule
Familie & Bindender P10-Endstand\\
\midrule
\endhead
Primkanal-/Projektor-No-Gos
& konkrete Rang-, Projektor-, Positivitäts- und Mangoldt-Abkürzungen gesperrt\\
Primfaser-Spektrum
& auditierter Transportgenerator kein diskreter HP-Endoperator\\
Feshbach
& finite Identität gültig; globaler Schatten-/Fredholm-Transfer offen\\
Determinante NEU-088--90
& $D_N(z)\to1$; nichttrivialer $C\xi$-Grenzwert in dieser Skalierung ausgeschlossen\\
Historischer Wert $e^{-\gamma^2/4}$
& SUPERSEDED\\
LFF/Rampe
& LFF $\Rightarrow$ Rampe bewiesen$_{\rm part}$; Umkehrung OPEN\\
Jacobi Startvektor
& unrenormierter Kanal kollabiert; Alternativen offen\\
Primeclock H1
& ungewichtet NO-GO; gewichteter Ersatz OPEN\\
Mellin/Prime-only
& historische Prime-only-Importe korrigiert / teilweise SUPERSEDED\\
Finite Part
& analytische Fortsetzungsdefinition ist kein operatorieller Grenzbeweis\\
P09 Koeffizienten
& mehrere kanonische Glätter-/1D-/Standard-SAYD-Routen NO-GO\\
P09 Rotation
& kanonischer skalarer Basislift: $t\Phi_0\ne C\Phi_0$\\
Objekt X global
& ausdrücklich OPEN; von P10 nicht widerlegt\\
\bottomrule
\end{longtable}
\normalsize

% ----------------------------------------------------------------
\section{Provenienz und Manuskriptstatus}
% ----------------------------------------------------------------

\begin{remark}[Bindende Quellen]
Bindende Primärquelle für die vollständige ID-Auflösung P10-N01--P10-N54 und
P10-O01--P10-O29 ist die final reconciliierte Pass-A-Matrix
\texttt{AUDIT-2026-08-09\_P10\_PassA\_Inventar\_NoGo\_Matrix\_P05-P09.md}.
Der Pass-A FINAL SEAL, der pfadgebundene Gegencheck und der Targeted-Reaudit
P06/P07 besitzen die im Markdown-SYN festgelegte Präzedenz. Die direkte
Markdown-Quelle ist
\texttt{papers/P10\_No-Go\_Theorems\_for\_Canonical\_Global\_Coupling.md}
im Status \emph{SYN FINAL AUDITED}.
\end{remark}

\begin{remark}[Transferstatus]
Diese LaTeX-Fassung ist ausschließlich als inhaltsgleicher Transfer des
geprüften Markdown-Endstands angelegt. Sie darf erst nach eigenständigem
LaTeX-Transferaudit als eingefroren gelten. Erst danach ist P11 ---
\emph{Global Coupling and the Object-X Candidate Geometry} --- prozedural
freizugeben.
\end{remark}

\end{document}
